\section{Discussion}

The skill label here is a property of the surgeon, constant across their trials. Under the permissive protocol every test surgeon also appears in training, so a model that recognises an individual can retrieve that individual's label without representing skill, and will be scored correct for it. Our results are consistent with this: the protocol alone accounts for at least a third of reported accuracy, and the features identify the operator far above chance. We have not shown the classifier internally relies on identity, and a reader may hold the weaker view that both results reflect some third property of individuals; our data do not separate those accounts. Either way, a figure obtained with the surgeon present in training does not estimate performance on a new surgeon, the only situation in which such a tool would be deployed. Any dataset predicting a person-level label from repeated measurements of the same people admits the same shortcut, though we measured it only here.

What survives is modest: one task of three beats a reference using no motion information, by roughly a quarter to a third of the per-task standard deviation, and there a single-feature regression on duration matches both models with heavily overlapping intervals. The zero-shot model's output correlates with duration at -0.81 to -0.91, so the convergence is not an artefact of fitting. We disclose that its rubric, written before any trial was scored but after we had seen which features carried signal, refers to completion time and pauses; this is a plausible route by which it came to track duration and we cannot exclude it.

For deployment the relevant finding is that confidence does not track correctness on unfamiliar operators. Both paths are overconfident and score worse than a constant forecast; a system confidently wrong on new surgeons is worse for triage than one uncertain everywhere, because the errors it passes through unreviewed are the ones it will not flag. Neither path's confidence identified deferrable trials, but the two nulls are not equal evidence: one is clean, one uninterpretable by construction, and both under-powered. A predictor with native uncertainty, through quantile or conformal methods, would remove the pairing defect and is the improvement a follow-up should adopt.

Three cheap changes would have caught this: report the held-out-surgeon protocol as primary with the paired difference beside it; anchor claims to an uninformative reference, since a positive correlation can accompany larger error than predicting the training mean; and publish a subject-identifiability diagnostic alongside a one-feature baseline. We note the tension in recommending protocol discipline from a study that substituted its own headline metric and added its central diagnostic post hoc; we have labelled both rather than presenting the final analysis as the planned one.

This is a single-dataset case study. Eight surgeons, seven on one task, make every interval crude, and those intervals express which surgeons were evaluated rather than refitting variability or the sampling of these eight. We examined one baseline and one commercial model at one pinned version with one rubric. We reproduced no published system, so our claims concern the protocol's structure and our own demonstration of it, not the validity of specific prior results. The confound and the calibration failure are demonstrated here, not established as general, and both should be checked on a second, larger dataset before being treated as properties of the field.
